\documentclass[11pt,a4paper]{article}
\usepackage[utf8]{inputenc}
\usepackage[margin=1in]{geometry}
\usepackage{amsmath}
\usepackage{amssymb}
\usepackage{enumitem}
\usepackage{booktabs}
\usepackage{titlesec}
\usepackage{microtype}

\titleformat{\section}{\large\bfseries\sffamily}{}{0em}{}[\titlerule]
\titleformat{\subsection}{\normalsize\bfseries\sffamily}{}{0em}{}

\setlength{\parindent}{0pt}
\setlength{\parskip}{6pt plus 2pt minus 1pt}

\title{\textbf{CAMS Exam: Select-Two Questions (50 Questions)}}
\author{\textbf{Advanced AML Training}}
\date{\today}

\begin{document}
	
	\maketitle
	
	\section*{Instructions}
	Each question has \textbf{five options (A through E)}. Exactly \textbf{two} of these options are correct. Select both correct answers.
	
	\section{Money Laundering Fundamentals}
	
	\subsection{1. Money Laundering Stages}
	Which two of the following correctly describe the three stages of money laundering?
	\begin{enumerate}[label=\Alph*.]
		\item Placement – introducing illicit funds into the financial system.
		\item Verification – confirming customer identity through biometrics.
		\item Layering – conducting complex transactions to obscure the audit trail.
		\item Approval – obtaining senior management sign-off for high-risk accounts.
		\item Integration – making laundered funds appear legitimate through investments.
	\end{enumerate}
	\textbf{Correct Answers: A and C}
	
	\subsection{2. Trade-Based Money Laundering (TBML) Indicators}
	Which two of the following are red flags for trade-based money laundering?
	\begin{enumerate}[label=\Alph*.]
		\item Invoice prices significantly above or below market value.
		\item Direct shipping between two FATF member countries with no intermediaries.
		\item Multiple amendments to letters of credit without logical explanation.
		\item Use of standard, unmodified commercial invoices.
		\item Goods shipped directly from manufacturer to end-user.
	\end{enumerate}
	\textbf{Correct Answers: A and C}
	
	\subsection{3. Shell Company Risk Indicators}
	Which two characteristics most increase money laundering risk for shell companies?
	\begin{enumerate}[label=\Alph*.]
		\item Physical office presence in a major financial center.
		\item Use of bearer shares or nominee shareholders to obscure ownership.
		\item Registration in a jurisdiction with public UBO registry.
		\item Multiple layers of corporate entities across secrecy havens.
		\item Single, clearly identified beneficial owner.
	\end{enumerate}
	\textbf{Correct Answers: B and D}
	
	\subsection{4. Black Market Peso Exchange (BMPE)}
	Which two statements accurately describe the Black Market Peso Exchange system?
	\begin{enumerate}[label=\Alph*.]
		\item It uses drug dollars in the U.S. to pay legitimate invoices for foreign merchants.
		\item It requires formal correspondent banking relationships to function.
		\item It allows Colombian peso brokers to profit from the exchange rate spread.
		\item It is a fully regulated currency exchange mechanism approved by FinCEN.
		\item It eliminates the need for any cash smuggling across borders.
	\end{enumerate}
	\textbf{Correct Answers: A and C}
	
	\subsection{5. Human Trafficking Financial Indicators}
	Which two transaction patterns most strongly indicate human trafficking for forced labor?
	\begin{enumerate}[label=\Alph*.]
		\item Single large wire transfer from a family member to a migrant worker.
		\item Multiple small, regular payments from a corporate account to multiple individuals with the same surname, no employment contracts.
		\item Frequent utility bill payments for multiple high-density residential properties from a single account.
		\item One-time purchase of international airline tickets for a family vacation.
		\item Large real estate purchase using verified inheritance funds.
	\end{enumerate}
	\textbf{Correct Answers: B and C}
	
	\section{FATF and International Standards}
	
	\subsection{6. FATF Mutual Evaluation Process}
	Which two statements about FATF Mutual Evaluations are correct?
	\begin{enumerate}[label=\Alph*.]
		\item Countries are assessed on both Technical Compliance and Effectiveness.
		\item Only developing countries are subject to mutual evaluations.
		\item Effectiveness is measured across 11 Immediate Outcomes.
		\item Technical Compliance is measured across 5 Recommendations only.
		\item Mutual evaluations are conducted annually for all member countries.
	\end{enumerate}
	\textbf{Correct Answers: A and C}
	
	\subsection{7. FATF Grey List vs. Black List}
	Which two statements accurately distinguish FATF Grey List and Black List jurisdictions?
	\begin{enumerate}[label=\Alph*.]
		\item Grey List jurisdictions are subject to increased monitoring but not full countermeasures.
		\item Black List jurisdictions are automatically expelled from the United Nations.
		\item Both lists require all financial transactions to be frozen immediately.
		\item Grey List jurisdictions have committed to resolve identified AML deficiencies within agreed timeframes.
		\item Black List jurisdictions receive simplified due diligence treatment to encourage compliance.
	\end{enumerate}
	\textbf{Correct Answers: A and D}
	
	\subsection{8. FATF Recommendation 16 (Wire Transfer Rule)}
	Which two pieces of information must accompany a cross-border wire transfer under FATF Recommendation 16?
	\begin{enumerate}[label=\Alph*.]
		\item Originator's full name.
		\item Originator's country of birth.
		\item Originator's account number.
		\item Originator's favorite color.
		\item Originator's mother's maiden name.
	\end{enumerate}
	\textbf{Correct Answers: A and C}
	
	\subsection{9. UN Security Council Sanctions}
	Which two statements correctly describe UNSC sanctions obligations for financial institutions?
	\begin{enumerate}[label=\Alph*.]
		\item Assets of designated individuals or entities must be frozen immediately without prior notice.
		\item Institutions may continue transactions if the amount is below \$10,000.
		\item The UN Consolidated Sanctions List is updated only once per year.
		\item Designated entities include terrorists, proliferators, and their financiers.
		\item Institutions may notify the customer before freezing assets to allow them to respond.
	\end{enumerate}
	\textbf{Correct Answers: A and D}
	
	\subsection{10. EU 5AMLD vs. 6AMLD}
	Which two features distinguish the EU Sixth Anti-Money Laundering Directive (6AMLD) from 5AMLD?
	\begin{enumerate}[label=\Alph*.]
		\item 6AMLD harmonizes predicate offenses across all EU member states.
		\item 6AMLD lowers the threshold for cash payments to €1,000.
		\item 6AMLD eliminates all customer due diligence requirements.
		\item 6AMLD introduces criminal liability for legal persons (corporations).
		\item 6AMLD removes all reporting obligations for virtual asset service providers.
	\end{enumerate}
	\textbf{Correct Answers: A and D}
	
	\section{USA PATRIOT Act and U.S. Regulations}
	
	\subsection{11. USA PATRIOT Act Section 311}
	Which two actions can the U.S. Secretary of the Treasury take under Section 311 against a foreign financial institution of primary money laundering concern?
	\begin{enumerate}[label=\Alph*.]
		\item Prohibit U.S. financial institutions from opening or maintaining correspondent accounts for the designated institution.
		\item Seize all assets of the foreign institution's U.S.-based executives personally.
		\item Impose special recordkeeping and reporting requirements on transactions involving the designated institution.
		\item Automatically freeze all individual customer accounts at the foreign institution.
		\item Mandate that all U.S. banks close their domestic retail branches.
	\end{enumerate}
	\textbf{Correct Answers: A and C}
	
	\subsection{12. Section 313: Shell Bank Prohibition}
	Which two statements accurately describe the U.S. prohibition on shell banks under Section 313 of the USA PATRIOT Act?
	\begin{enumerate}[label=\Alph*.]
		\item A shell bank is defined as a bank with no physical presence in any country.
		\item Shell banks are permitted if they are regulated by a G20 central bank.
		\item U.S. banks are prohibited from maintaining correspondent accounts for foreign shell banks.
		\item Shell banks may operate freely if they hold a valid state banking license.
		\item An exception exists for shell banks that are affiliates of U.S.-regulated banks.
	\end{enumerate}
	\textbf{Correct Answers: A and C}
	
	\subsection{13. SAR Safe Harbor Provisions}
	Which two protections does the SAR Safe Harbor provide to financial institutions that file Suspicious Activity Reports in good faith?
	\begin{enumerate}[label=\Alph*.]
		\item Immunity from civil liability arising from the SAR filing.
		\item Automatic approval for all future transactions involving the same customer.
		\item Protection from criminal liability for the underlying suspicious activity.
		\item Immunity from civil liability for failure to disclose the SAR to the customer.
		\item Guaranteed confidentiality of the SAR filing even if disclosed to the customer.
	\end{enumerate}
	\textbf{Correct Answers: A and D}
	
	\subsection{14. FinCEN Geographic Targeting Orders (GTOs)}
	Which two statements accurately describe FinCEN Geographic Targeting Orders for real estate transactions?
	\begin{enumerate}[label=\Alph*.]
		\item GTOs require title insurance companies to identify beneficial owners behind all-cash purchases above a threshold.
		\item GTOs apply to all real estate transactions regardless of purchase method.
		\item GTOs are permanent regulations that do not expire.
		\item GTOs target specific metropolitan areas identified as high-risk for real estate money laundering.
		\item GTOs exempt all transactions involving publicly traded companies.
	\end{enumerate}
	\textbf{Correct Answers: A and D}
	
	\subsection{15. FinCEN 314(a) vs. 314(b)}
	Which two statements correctly distinguish FinCEN 314(a) from 314(b)?
	\begin{enumerate}[label=\Alph*.]
		\item 314(a) involves mandatory searches of bank records in response to law enforcement requests.
		\item 314(b) involves voluntary information sharing among financial institutions.
		\item 314(a) requires customer notification before any search is conducted.
		\item 314(b) prohibits any information sharing under all circumstances.
		\item Both 314(a) and 314(b) require a court order to initiate.
	\end{enumerate}
	\textbf{Correct Answers: A and B}
	
	\section{CDD, KYC, and Beneficial Ownership}
	
	\subsection{16. Customer Due Diligence (CDD) Pillars}
	Under FinCEN's CDD Rule, which two of the following are required elements of CDD?
	\begin{enumerate}[label=\Alph*.]
		\item Identify and verify the identity of each customer.
		\item Obtain a credit score for every customer over \$1 million.
		\item Identify and verify beneficial owners of legal entity customers (25\% ownership or control prong).
		\item Conduct annual physical inspection of all business premises.
		\item Monitor transactions on an ongoing basis for suspicious activity.
	\end{enumerate}
	\textbf{Correct Answers: A and C}
	
	\subsection{17. Beneficial Ownership Threshold Calculation}
	Which two scenarios would require an individual to be identified as a UBO under a standard 25\% ownership threshold?
	\begin{enumerate}[label=\Alph*.]
		\item Individual A directly owns 20\% of Target Company.
		\item Individual B directly owns 15\% and indirectly owns another 12\% through a wholly-owned holding company.
		\item Individual C directly owns 24.9\% of Target Company.
		\item Individual D owns 30\% of Parent Company, which owns 100\% of Target Company.
		\item Individual E owns 10\% of Target Company and serves as a junior administrative assistant.
	\end{enumerate}
	\textbf{Correct Answers: B and D}
	
	\subsection{18. Politically Exposed Persons (PEPs)}
	Which two individuals would be classified as PEPs or PEP family members under international standards?
	\begin{enumerate}[label=\Alph*.]
		\item A foreign head of state.
		\item A senior cabinet minister in a foreign government.
		\item A local grocery store cashier in a low-risk jurisdiction.
		\item The adult child of a foreign head of state.
		\item A retired elementary school teacher with no government connections.
	\end{enumerate}
	\textbf{Correct Answers: A and D}
	
	\subsection{19. Enhanced Due Diligence (EDD) Requirements}
	Which two actions are required as part of Enhanced Due Diligence for high-risk customers?
	\begin{enumerate}[label=\Alph*.]
		\item Establishing the source of wealth and source of funds through verifiable documentation.
		\item Applying simplified identity verification (e.g., selfie only).
		\item Obtaining senior management approval before onboarding or maintaining the relationship.
		\item Waiving all transaction monitoring for privacy reasons.
		\item Conducting a credit check regardless of the product type.
	\end{enumerate}
	\textbf{Correct Answers: A and C}
	
	\subsection{20. Simplified Due Diligence (SDD)}
	Which two types of customers may qualify for Simplified Due Diligence under a risk-based approach?
	\begin{enumerate}[label=\Alph*.]
		\item A publicly traded company listed on a recognized stock exchange with full SEC disclosure.
		\item A foreign PEP from a high-corruption jurisdiction.
		\item A U.S. government agency (e.g., Social Security Administration).
		\item A shell company registered in a secrecy haven.
		\item A customer who refuses to provide any identification documents.
	\end{enumerate}
	\textbf{Correct Answers: A and C}
	
	\section{Virtual Assets and VASPs}
	
	\subsection{21. FATF Travel Rule for Virtual Assets}
	Under the FATF Travel Rule for virtual assets, which two pieces of information must VASPs collect and share for transactions above the threshold?
	\begin{enumerate}[label=\Alph*.]
		\item Originator's name and account number (wallet identifier).
		\item Originator's favorite cryptocurrency trading strategy.
		\item Beneficiary's name and account number (wallet identifier).
		\item The exact GPS coordinates of the transaction initiation.
		\item The originator's internet browsing history.
	\end{enumerate}
	\textbf{Correct Answers: A and C}
	
	\subsection{22. Anonymity-Enhanced Cryptocurrencies (AECs)}
	Which two characteristics of Monero (XMR) and similar privacy coins create money laundering risks?
	\begin{enumerate}[label=\Alph*.]
		\item They use ring signatures and stealth addresses to obscure transaction trails.
		\item They are fully transparent and traceable on public blockchains.
		\item They are illegal in all jurisdictions globally.
		\item They can be used with mixers/tumblers to further obfuscate the audit trail.
		\item They are exclusively issued and controlled by central banks.
	\end{enumerate}
	\textbf{Correct Answers: A and D}
	
	\subsection{23. Crypto Mixers and Tumblers}
	Which two risk indicators should trigger enhanced scrutiny when a customer uses a crypto mixing service?
	\begin{enumerate}[label=\Alph*.]
		\item The customer deposits funds into a mixer and withdraws to a new wallet.
		\item The customer provides a clear, documented legitimate reason for using the mixer.
		\item The customer refuses to explain the source of funds before mixing.
		\item The customer uses only regulated, licensed mixers with full KYC.
		\item The customer immediately sends mixed funds to a known illicit address.
	\end{enumerate}
	\textbf{Correct Answers: A and C}
	
	\subsection{24. Decentralized Finance (DeFi) AML Responsibility}
	Under FATF guidance, which two entities may bear AML/CFT obligations in DeFi arrangements?
	\begin{enumerate}[label=\Alph*.]
		\item The individual developers who retain control or governance over the protocol.
		\item Every anonymous user of the protocol equally.
		\item The blockchain network itself as a legal person.
		\item Any entity that exercises control or sufficient influence over the protocol's operations.
		\item No entity ever bears obligations because DeFi is by definition unregulated.
	\end{enumerate}
	\textbf{Correct Answers: A and D}
	
	\subsection{25. Stablecoin Money Laundering Risks}
	Which two characteristics of stablecoins create money laundering vulnerabilities?
	\begin{enumerate}[label=\Alph*.]
		\item They are pegged to fiat currency, providing price stability that facilitates value storage.
		\item They are always fully regulated by central banks with complete transaction transparency.
		\item They can be transferred rapidly across borders without traditional banking intermediaries.
		\item They cannot be used in any DeFi protocol or decentralized exchange.
		\item Some stablecoin issuers have weak or no CDD controls on their direct customers.
	\end{enumerate}
	\textbf{Correct Answers: A and C}
	
	\section{Correspondent Banking and Cross-Border Risks}
	
	\subsection{26. Nested Correspondent Accounts}
	Which two risks arise when a respondent bank allows third-party banks to route nested transactions through a correspondent account without disclosure?
	\begin{enumerate}[label=\Alph*.]
		\item The correspondent bank becomes completely blind to the true identity of the ultimate originators.
		\item The risk of processing transactions for sanctioned entities without knowledge increases significantly.
		\item Operational costs decrease because fewer transactions need review.
		\item The respondent bank's compliance burden is eliminated entirely.
		\item The correspondent bank gains perfect visibility into every transaction.
	\end{enumerate}
	\textbf{Correct Answers: A and B}
	
	\subsection{27. Payable-Through Accounts (PTAs)}
	Which two obligations apply to U.S. correspondent banks that offer payable-through accounts to respondent banks?
	\begin{enumerate}[label=\Alph*.]
		\item The correspondent bank must apply CDD requirements directly to the respondent bank's underlying customers.
		\item The correspondent bank may ignore all transactions processed through the PTA.
		\item The correspondent bank must maintain records identifying each PTA accountholder.
		\item The correspondent bank must allow unlimited anonymous transactions.
		\item The correspondent bank must outsource all monitoring to the respondent bank without oversight.
	\end{enumerate}
	\textbf{Correct Answers: A and C}
	
	\subsection{28. De-Risking Consequences}
	Which two negative consequences result from indiscriminate de-risking by major correspondent banks?
	\begin{enumerate}[label=\Alph*.]
		\item Legitimate remittance corridors in developing economies lose access to global finance.
		\item The global financial system becomes more transparent and inclusive.
		\item Funds are driven into unregulated, informal channels that are harder to monitor.
		\item Money laundering risks are completely eliminated from the financial system.
		\item Regulatory authorities universally praise de-risking as best practice.
	\end{enumerate}
	\textbf{Correct Answers: A and C}
	
	\subsection{29. Wolfsberg Correspondent Banking Due Diligence Questionnaire (CDDQ)}
	Which two areas are typically covered in a Wolfsberg Correspondent Banking Due Diligence Questionnaire?
	\begin{enumerate}[label=\Alph*.]
		\item The respondent bank's AML/CFT policies, procedures, and internal controls.
		\item The respondent bank's CEO's personal investment portfolio.
		\item The respondent bank's customer base, including high-risk segments like MSBs and NBFIs.
		\item The respondent bank's office furniture and decoration preferences.
		\item The respondent bank's employee dress code policy.
	\end{enumerate}
	\textbf{Correct Answers: A and C}
	
	\subsection{30. Reverse Correspondent Banking}
	Which two risks are unique to reverse correspondent banking arrangements (small bank providing account to larger international bank)?
	\begin{enumerate}[label=\Alph*.]
		\item The small bank may be used as a blind gateway into the larger bank's client base.
		\item The small bank gains perfect visibility into all of the larger bank's global transactions.
		\item The small bank may struggle to conduct adequate due diligence on the larger bank's complex customer structures.
		\item The reverse arrangement eliminates all money laundering risk for both banks.
		\item The large bank automatically assumes full liability for all small bank transactions.
	\end{enumerate}
	\textbf{Correct Answers: A and C}
	
	\section{Proliferation Financing and Sanctions}
	
	\subsection{31. Proliferation Financing Red Flags}
	Which two indicators suggest potential proliferation financing (weapons of mass destruction-related) activity?
	\begin{enumerate}[label=\Alph*.]
		\item Export of dual-use goods to a country under UNSC nuclear sanctions.
		\item Payment routed through third-party intermediaries unrelated to the transaction.
		\item Shipment of agricultural products between two FATF-compliant nations.
		\item Transaction conducted entirely in local currency within a single country.
		\item End-user listed as a legitimate, publicly-traded pharmaceutical company.
	\end{enumerate}
	\textbf{Correct Answers: A and B}
	
	\subsection{32. UNSCR 1540 Obligations}
	Which two actions are required of UN member states under UN Security Council Resolution 1540?
	\begin{enumerate}[label=\Alph*.]
		\item Establish effective controls over chemical, biological, radiological, and nuclear (CBRN) materials.
		\item Freeze all assets of any country developing any form of advanced technology.
		\item Prohibit any non-state actor from manufacturing or acquiring WMD-related materials.
		\item Eliminate all border controls between member states.
		\item Require all private citizens to report all scientific research to the UN.
	\end{enumerate}
	\textbf{Correct Answers: A and C}
	
	\subsection{33. OFAC 50\% Rule}
	Which two statements correctly describe the OFAC 50\% Rule for blocked persons?
	\begin{enumerate}[label=\Alph*.]
		\item Any entity owned 50\% or more in the aggregate by one or more blocked persons is itself blocked.
		\item Ownership is calculated by multiplying through each tier of ownership.
		\item Only direct ownership matters; indirect ownership is ignored entirely.
		\item The rule applies only to individuals, not to corporate entities.
		\item The 50\% threshold is applied regardless of whether the ownership is direct or indirect.
	\end{enumerate}
	\textbf{Correct Answers: A and B}
	
	\subsection{34. Secondary Sanctions}
	Which two statements accurately describe secondary sanctions risk for non-U.S. financial institutions?
	\begin{enumerate}[label=\Alph*.]
		\item A non-U.S. bank that processes non-USD transactions with a sanctioned Iranian entity may face U.S. secondary sanctions.
		\item Secondary sanctions can result in the non-U.S. bank being designated on the OFAC SDN list.
		\item Secondary sanctions never apply to any transaction that avoids the U.S. dollar.
		\item The EU Blocking Statute completely eliminates all secondary sanctions risk for European banks.
		\item Secondary sanctions may restrict a non-U.S. bank's access to the U.S. financial system.
	\end{enumerate}
	\textbf{Correct Answers: A and E}
	
	\subsection{35. General vs. Specific OFAC Licenses}
	Which two statements correctly distinguish between OFAC General Licenses and Specific Licenses?
	\begin{enumerate}[label=\Alph*.]
		\item A General License authorizes certain categories of transactions without the need to apply to OFAC.
		\item A Specific License is issued to a particular person or entity for a specific transaction.
		\item General Licenses must be renewed annually with a formal application to OFAC.
		\item Specific Licenses are published publicly and apply to all persons globally without application.
		\item Both license types require a \$10,000 application fee paid to the U.S. Treasury.
	\end{enumerate}
	\textbf{Correct Answers: A and B}
	
	\section{AML Program Governance and Testing}
	
	\subsection{36. Three Lines of Defense Model}
	Which two functions correctly align with the Three Lines of Defense in AML governance?
	\begin{enumerate}[label=\Alph*.]
		\item First Line: Business operations and front-line staff who own customer relationships.
		\item Second Line: Independent Audit Department that designs transaction monitoring rules.
		\item Second Line: Compliance Department and AML Officer providing oversight and guidance.
		\item Third Line: Business line executives who approve every transaction.
		\item Third Line: Internal Audit providing independent assurance to the Board.
	\end{enumerate}
	\textbf{Correct Answers: A and C}
	
	\subsection{37. Independent Testing Requirements}
	Which two characteristics are required for the independent testing (audit) of an AML program?
	\begin{enumerate}[label=\Alph*.]
		\item The auditor must be independent of the functions being tested.
		\item The auditor must be the same person who designed the AML policies.
		\item The auditor must have sufficient expertise and qualifications.
		\item The auditor must report directly to the Chief Compliance Officer only.
		\item The auditor must be employed by the same business line being audited.
	\end{enumerate}
	\textbf{Correct Answers: A and C}
	
	\subsection{38. Board of Directors AML Responsibilities}
	Which two responsibilities belong directly to a financial institution's Board of Directors?
	\begin{enumerate}[label=\Alph*.]
		\item Approving the AML compliance program and its material changes.
		\item Filing every Suspicious Activity Report individually.
		\item Reviewing every customer transaction over \$1,000.
		\item Providing oversight and holding senior management accountable for AML effectiveness.
		\item Conducting daily branch teller transactions.
	\end{enumerate}
	\textbf{Correct Answers: A and D}
	
	\subsection{39. Transaction Monitoring False Negatives vs. False Positives}
	Which two statements correctly characterize AML transaction monitoring calibration risks?
	\begin{enumerate}[label=\Alph*.]
		\item False positives occur when legitimate activity triggers an alert that is not suspicious.
		\item False negatives occur when suspicious activity fails to generate any alert.
		\item False positives are always preferable to false negatives because they cause no harm.
		\item False negatives are acceptable as long as the bank files at least one SAR per year.
		\item Over-filtering (raising thresholds too high) increases false negatives.
	\end{enumerate}
	\textbf{Correct Answers: A and B}
	
	\subsection{40. AML Training Requirements}
	Which two statements correctly describe effective AML training programs under regulatory expectations?
	\begin{enumerate}[label=\Alph*.]
		\item Training must be role-specific and tailored to different employee functions.
		\item All employees must receive identical training regardless of role.
		\item Board members and senior management require specialized training on governance and risk oversight.
		\item Training is required only once at employee onboarding with no refreshers needed.
		\item Frontline tellers do not need any AML training because they do not file SARs.
	\end{enumerate}
	\textbf{Correct Answers: A and C}
	
	\section{SAR Filing, Tipping-Off, and Investigations}
	
	\subsection{41. SAR Filing Confidentiality (Tipping-Off)}
	Which two actions would constitute illegal ``tipping-off'' under AML laws?
	\begin{enumerate}[label=\Alph*.]
		\item Telling a customer that a SAR has been filed on their account.
		\item Disclosing to a third party that an investigation is being considered.
		\item Filing a SAR in good faith with the relevant FIU.
		\item Sharing a SAR internally with the compliance committee on a need-to-know basis.
		\item Publicly announcing on social media that a specific customer is under investigation.
	\end{enumerate}
	\textbf{Correct Answers: A and E}
	
	\subsection{42. Responding to Law Enforcement Requests}
	Which two actions are appropriate when a financial institution receives a valid law enforcement subpoena for customer records?
	\begin{enumerate}[label=\Alph*.]
		\item Promptly comply with the subpoena within the specified timeframe.
		\item Notify the customer immediately so they can close the account.
		\item Maintain a complete internal record of all documents produced.
		\item Destroy all responsive records to protect customer privacy.
		\item Refuse to comply unless the law enforcement officer provides a verbal promise of confidentiality.
	\end{enumerate}
	\textbf{Correct Answers: A and C}
	
	\subsection{43. Law Enforcement ``Keep Open'' Requests}
	Which two protocols should a bank follow when law enforcement requests that a suspicious account remain open for ongoing surveillance?
	\begin{enumerate}[label=\Alph*.]
		\item Document the request in writing from law enforcement.
		\item Close the account immediately regardless of the request.
		\item Obtain senior management approval before agreeing to the request.
		\item Publicly announce the surveillance operation to deter criminal activity.
		\item Continue normal transaction monitoring without any enhanced scrutiny.
	\end{enumerate}
	\textbf{Correct Answers: A and C}
	
	\subsection{44. Mutual Legal Assistance Treaties (MLATs)}
	Which two statements correctly describe the function of Mutual Legal Assistance Treaties?
	\begin{enumerate}[label=\Alph*.]
		\item MLATs provide a formal mechanism for cross-border exchange of evidence and judicial documents.
		\item MLATs allow one country to unilaterally arrest citizens of another country without authorization.
		\item MLATs facilitate asset freezing and forfeiture across international borders.
		\item MLATs eliminate all data privacy laws in both signatory countries.
		\item MLATs are informal, non-binding agreements with no legal force.
	\end{enumerate}
	\textbf{Correct Answers: A and C}
	
	\subsection{45. Egmont Group FIU-to-FIU Exchange}
	Which two principles govern information exchange between Financial Intelligence Units under the Egmont Group framework?
	\begin{enumerate}[label=\Alph*.]
		\item Exchanged information must be used only for AML/CFT purposes.
		\item Information may be freely shared with the subject of the intelligence.
		\item The receiving FIU must maintain confidentiality of the shared intelligence.
		\item No limitations exist on how the receiving FIU may use the information.
		\item Information may be sold to commercial data brokers for profit.
	\end{enumerate}
	\textbf{Correct Answers: A and C}
	
	\section{MSBs, Casinos, and Gatekeepers}
	
	\subsection{46. Money Services Business (MSB) Registration}
	Which two activities trigger MSB registration requirements under U.S. federal law?
	\begin{enumerate}[label=\Alph*.]
		\item Check cashing exceeding \$1,000 for any person in a single day.
		\item Operating a federally insured commercial bank.
		\item Money transmitting services regardless of volume.
		\item Selling homemade crafts at a local flea market.
		\item Issuing or selling prepaid access cards.
	\end{enumerate}
	\textbf{Correct Answers: A and C}
	
	\subsection{47. Casino AML Obligations}
	Under the U.S. Bank Secrecy Act, which two reporting obligations apply to casinos with gross gaming revenue exceeding \$1 million?
	\begin{enumerate}[label=\Alph*.]
		\item Currency Transaction Reports (CTRs) for cash transactions exceeding \$10,000 in a gaming day.
		\item Suspicious Activity Reports (SARs) for transactions exceeding \$5,000 that exhibit suspicious patterns.
		\item Complete exemption from all BSA reporting requirements.
		\item Daily reports of all winnings regardless of amount.
		\item Reports only on slot machine activity, not table games.
	\end{enumerate}
	\textbf{Correct Answers: A and B}
	
	\subsection{48. Lawyer Trust Accounts (IOLTA) Vulnerabilities}
	Which two characteristics make IOLTA accounts vulnerable to money laundering exploitation?
	\begin{enumerate}[label=\Alph*.]
		\item Banks often rely on the attorney's professional status without looking through to individual client sub-transactions.
		\item IOLTA accounts are completely immune from any law enforcement subpoena.
		\item Multiple client funds are pooled together, obscuring the origin of specific funds.
		\item IOLTA accounts are required to be fully transparent to the public.
		\item Lawyers are never required to identify their clients for any purpose.
	\end{enumerate}
	\textbf{Correct Answers: A and C}
	
	\subsection{49. Informal Value Transfer Systems (IVTS) / Hawala}
	Which two characteristics define hawala and similar IVTS arrangements?
	\begin{enumerate}[label=\Alph*.]
		\item Transactions are settled through trust-based obligations rather than formal banking records.
		\item Every hawala transaction is recorded on a public blockchain.
		\item Hawaladars (brokers) maintain informal ledgers to track reciprocal obligations.
		\item IVTS systems are always illegal in all countries without exception.
		\item Hawala cannot be used for amounts exceeding \$100.
	\end{enumerate}
	\textbf{Correct Answers: A and C}
	
	\subsection{50. FATF Recommendation 8: NPOs}
	Which two vulnerabilities make certain non-profit organizations susceptible to terrorist financing exploitation?
	\begin{enumerate}[label=\Alph*.]
		\item NPOs often operate in conflict zones where financial monitoring infrastructure is weak.
		\item NPOs are always fully regulated by central banks with complete transaction oversight.
		\item NPOs enjoy public trust and may receive large cash donations without rigorous scrutiny.
		\item NPOs are legally prohibited from having bank accounts in any jurisdiction.
		\item All NPOs are required to publish their donor lists publicly.
	\end{enumerate}
	\textbf{Correct Answers: A and C}
	
\end{document}